\heading{数学物理方法（下）-模拟题1-期中}

\noteinfo{题目和答案由AI生成。}

\begin{question}
设一根长度为 $l$ 的均匀细杆作纵振动。杆受到与速度成正比的阻尼和与位移成正比的回复力，位移记为 $u(x,t)$，其中 $0<x<l$。杆的上端固定，下端自由。
\begin{enumerate}[label=（\arabic*）,leftmargin=2.8em,itemsep=0.2em]
\item 写出控制方程与边界条件；
\item 令 $u(x,t)=X(x)T(t)$，写出相应的分离后常微分方程；
\item 写出空间本征值问题，并给出本征值、本征函数与模方。
\end{enumerate}
\begin{solution}

\begin{enumerate}[label=（\arabic*）,leftmargin=2.8em,itemsep=0.2em]
\item 可写为
\[
\rho u_{tt}+d u_t-Eu_{xx}+ku=0,
\qquad 0<x<l,
\]
其中 $\rho$ 为线密度，$E$ 为等效弹性常数。上端固定、下端自由对应
\[
 u(0,t)=0,\qquad u_x(l,t)=0.
\]
\item 代入 $u=XT$ 并除以 $XT$，得
\[
\rho \frac{T''}{T}+d\frac{T'}{T}+k=E\frac{X''}{X}=-\lambda.
\]
故
\[
X''+\frac{\lambda}{E}X=0,
\qquad
\rho T''+dT'+(k+\lambda)T=0.
\]
\item 空间本征值问题为
\[
X''+\mu^2X=0,
\qquad X(0)=0,\quad X'(l)=0,
\qquad \mu^2=\frac{\lambda}{E}.
\]
通解 $X=A\cos\mu x+B\sin\mu x$。由 $X(0)=0$ 得 $A=0$，再由 $X'(l)=0$ 得
\[
\cos(\mu l)=0\Longrightarrow \mu_n=\frac{(n+\tfrac12)\pi}{l},\quad n=0,1,2,\dots
\]
故
\[
\lambda_n=E\mu_n^2=E\Bigl(\frac{(n+\tfrac12)\pi}{l}\Bigr)^2,
\qquad X_n(x)=\sin\frac{(n+\tfrac12)\pi x}{l}.
\]
模方为
\[
\|X_n\|^2=\int_0^l X_n(x)^2\,dx=\frac l2.
\]
归一化后可取 $\sqrt{2/l}\,X_n(x)$。
\end{enumerate}
\ansmath{\lambda_n=E\mu_n^2=E\Bigl(\frac{(n+\tfrac12)\pi}{l}\Bigr)^2,
\qquad X_n(x)=\sin\frac{(n+\tfrac12)\pi x}{l}\quad \|X_n\|^2=\int_0^l X_n(x)^2\,dx=\frac l2}

\end{solution}
\end{question}

\begin{question}
讨论下列势函数对应的薛定谔方程在哪些坐标系中可分离变量，并简述理由：
\[
V_1(x,y,z)=\alpha(x^2+y^2+z^2),\qquad V_2(x,y,z)=\beta z^2,
\qquad V_3(x,y,z)=\frac{\gamma}{x^2+y^2+z^2}.
\]
可选坐标系：直角坐标、柱坐标、球坐标。
\begin{solution}

\begin{enumerate}[label=（\arabic*）,leftmargin=2.8em,itemsep=0.2em]
\item $V_1=\alpha r^2$ 是中心势，因此在球坐标可分离；又因 $r^2=x^2+y^2+z^2$ 是三个单变量项之和，所以在直角坐标也可分离；在柱坐标中 $r^2=\rho^2+z^2$，同样可分离。故三种坐标都可分离。
\item $V_2=\beta z^2$ 只依赖 $z$，因此在直角坐标中可分离；在柱坐标中仍只依赖 $z$，也可分离；在球坐标中 $z=r\cos\theta$，势写成 $\beta r^2\cos^2\theta$，耦合了 $r,\theta$，一般不可分离。
\item $V_3=\gamma/r^2$ 仅依赖 $r$，在球坐标可分离；在柱坐标中写成 $\gamma/(\rho^2+z^2)$，耦合 $\rho,z$，不可分离；在直角坐标中写成 $\gamma/(x^2+y^2+z^2)$，耦合三个变量，也不可分离。
\end{enumerate}
\anstext{见上。}

\end{solution}
\end{question}

\begin{question}
考虑区间 $0<x<l$ 上的算符
\[
Ly=-y'',
\]
边界条件为
\[
y(l)=\alpha y(0),\qquad y'(l)=\beta y'(0),
\]
其中 $\alpha,\beta\in\mathbb C$ 为常数。
\begin{enumerate}[label=（\arabic*）,leftmargin=2.8em,itemsep=0.2em]
\item 求该边值问题自伴的条件；
\item 在自伴条件下，说明本征值为实数；
\item 当 $\alpha=\beta=1$ 时，求全部本征值和本征函数。
\end{enumerate}
\begin{solution}

对任意足够光滑的 $y,z$，Green 公式给出
\[
\langle Ly,z\rangle-\langle y,Lz\rangle=[-y'\bar z+y\bar z']_0^l.
\]
由边界条件代入，边界项化为
\[
-y'(l)\overline{z(l)}+y(l)\overline{z'(l)}+y'(0)\overline{z(0)}-y(0)\overline{z'(0)}
=(1-\beta\overline\alpha)y'(0)\overline{z(0)}+(\alpha\overline\beta-1)y(0)\overline{z'(0)}.
\]
因此对所有满足同类边界条件的 $y,z$ 都为零，当且仅当
\[
\alpha\overline\beta=1.
\]
这就是自伴条件。

自伴后，标准 Sturm--Liouville 理论立即推出本征值为实数；也可令 $Ly=\lambda y$，取内积得
\[
\lambda\langle y,y\rangle=\langle Ly,y\rangle\in\mathbb R.
\]

当 $\alpha=\beta=1$ 时是周期边界条件。方程
\[
-y''=\lambda y
\]
分三种情形：
\begin{enumerate}[label=（\arabic*）,leftmargin=2.8em,itemsep=0.2em]
\item $\lambda=0$：解为常数，故有本征值 $\lambda_0=0$，本征函数 $y_0(x)=1$；
\item $\lambda=\mu^2>0$：通解 $A\cos\mu x+B\sin\mu x$。由周期条件得到
\[
\mu l=2n\pi,\qquad n=1,2,\dots
\]
故
\[
\lambda_n=\Bigl(\frac{2n\pi}{l}\Bigr)^2,
\quad y_n^{(1)}=\cos\frac{2n\pi x}{l},
\quad y_n^{(2)}=\sin\frac{2n\pi x}{l};
\]
\item $\lambda<0$ 时只有零解，不给本征值。
\end{enumerate}
\ansmath{\mu l=2n\pi,\qquad n=1,2,\dots\quad \lambda_n=\Bigl(\frac{2n\pi}{l}\Bigr)^2,
\quad y_n^{(1)}=\cos\frac{2n\pi x}{l},
\quad y_n^{(2)}=\sin\frac{2n\pi x}{l};}

\end{solution}
\end{question}

\begin{question}
求解定解问题
\[
\begin{cases}
 u_{tt}-a^2u_{xx}=0, & 0<x<\pi,\ t>0,\\
 u_x(0,t)=0,\quad u_x(\pi,t)=0, & t>0,\\
 u(x,0)=x(\pi-x),\quad u_t(x,0)=\cos 2x.
\end{cases}
\]

\begin{solution}

Neumann 边界对应本征函数系 $\{\cos nx\}_{n\ge0}$，故解设为
\[
u(x,t)=A_0+B_0t+\sum_{n=1}^{\infty}\bigl(A_n\cos nat+B_n\sin nat\bigr)\cos nx.
\]
先把初位移与初速度作余弦展开。

对 $f(x)=x(\pi-x)$，
\[
A_0=\frac1\pi\int_0^{\pi}x(\pi-x)\,dx=\frac{\pi^2}{6}.
\]
对 $n\ge1$，
\[
A_n=\frac{2}{\pi}\int_0^{\pi}x(\pi-x)\cos nx\,dx
=\begin{cases}
-\dfrac{4}{n^2},& n\text{ 偶},\\[0.3em]
0,& n\text{ 奇}.
\end{cases}
\]
又 $g(x)=\cos2x$ 说明只有 $n=2$ 的速度模态非零，因此
\[
B_2\cdot 2a=1\Longrightarrow B_2=\frac1{2a},\qquad B_n=0\ (n\ne2),\quad B_0=0.
\]
故
\ansmath{u(x,t)=\frac{\pi^2}{6}-\sum_{m=1}^{\infty}\frac1{m^2}\cos(2ma t)\cos(2mx)+\frac{1}{2a}\sin(2at)\cos 2x.}
其中用了 $A_{2m}=-1/m^2$。
\end{solution}
\end{question}

\begin{question}
求解热传导问题
\[
\begin{cases}
 u_t=\kappa u_{xx}, & 0<x<l,\ t>0,\\
 -K u_x(0,t)=q,\qquad u(l,t)=0, & t>0,\\
 u(x,0)=0.
\end{cases}
\]
要求写成“稳态项 $+$ 瞬态级数”的形式。
\begin{solution}

先求稳态解 $v(x)$：
\[
\kappa v''=0,
\qquad -K v'(0)=q,
\qquad v(l)=0.
\]
故
\[
 v(x)=\frac{q}{K}(l-x).
\]
令 $w=u-v$，则 $w$ 满足
\[
\begin{cases}
 w_t=\kappa w_{xx},\\
 w_x(0,t)=0,\quad w(l,t)=0,\\
 w(x,0)=-\dfrac{q}{K}(l-x).
\end{cases}
\]
相应本征问题
\[
X''+\mu^2X=0,
\qquad X'(0)=0,
\quad X(l)=0
\]
给出
\[
\mu_n=\frac{(n+\tfrac12)\pi}{l},\qquad X_n(x)=\cos\mu_n x.
\]
故
\[
 w(x,t)=\sum_{n=0}^{\infty}c_n e^{-\kappa\mu_n^2 t}\cos\mu_n x.
\]
系数为
\[
 c_n=\frac{2}{l}\int_0^l\left(-\frac{q}{K}(l-x)\right)\cos\mu_n x\,dx
 =-\frac{2q}{Kl\mu_n^2}.
\]
因此
\ansmath{u(x,t)=\frac{q}{K}(l-x)-\sum_{n=0}^{\infty}\frac{2q}{Kl\mu_n^2}e^{-\kappa\mu_n^2 t}\cos\mu_n x,
\qquad \mu_n=\frac{(n+\tfrac12)\pi}{l}.}
\end{solution}
\end{question}

\begin{question}
在上半环域
\[
 a<r<b,\qquad 0<\theta<\pi
\]
中求解 Laplace 方程
\[
\nabla^2u=0,
\]
满足边界条件
\[
\frac{\partial u}{\partial r}(a,\theta)=0,
\qquad u(b,\theta)=0,
\qquad u(r,0)=u(r,\pi)=0.
\]
写出分离变量后的本征结构，不必求出所有展开系数。
\begin{solution}

取 $u(r,\theta)=R(r)\Theta(\theta)$。代入后得
\[
\frac{r^2R''+rR'}{R}=-\frac{\Theta''}{\Theta}=\lambda.
\]
由于 $u(r,0)=u(r,\pi)=0$，角向满足
\[
\Theta''+n^2\Theta=0,
\qquad \Theta(0)=\Theta(\pi)=0,
\]
故
\[
\Theta_n(\theta)=\sin n\theta,
\qquad n=1,2,\dots
\]
对于每个 $n$，径向满足 Euler 方程
\[
r^2R''+rR'-n^2R=0,
\]
故
\[
R_n(r)=A_n r^n+B_n r^{-n}.
\]
边界条件
\[
R_n'(a)=0,
\qquad R_n(b)=0
\]
给出
\[
A_n a^{n-1}-B_n a^{-n-1}=0,
\qquad A_n b^n+B_n b^{-n}=0.
\]
解得比例关系
\[
B_n=nA_n a^{2n}
\quad\text{（或直接由两式联立给出）},
\]
更方便的写法是取
\[
R_n(r)=r^n+\frac{a^{2n}}{r^n},
\]
再用 $R_n(b)=0$ 调整整体常数，或写成
\[
R_n(r)=\left(r^n+\frac{a^{2n}}{r^n}\right)-\frac{b^n+a^{2n}b^{-n}}{b^n-a^{2n}b^{-n}}\left(r^n-\frac{a^{2n}}{r^n}\right).
\]
因此通解结构为
\ansmath{u(r,\theta)=\sum_{n=1}^{\infty}C_n R_n(r)\sin n\theta,}
其中 $R_n$ 满足上面的两个径向边界条件，系数由额外边界数据决定。
\end{solution}
\end{question}


